\section{Method: PA-DCT}
\label{sec:method}

PA-DCT is an instance of contextual Thompson sampling adapted to the agent-native regime (\S\ref{sec:intro}). Each (arm, task-type-bucket) cell maintains \emph{two} independent Gaussian posteriors---one over utility, one over cost---both tracked through discounted sufficient statistics. At decision time, the policy samples both posteriors and ranks affordable arms by a budget-pressure-weighted score; on receipt of the realized outcome, only the chosen cell's posteriors are updated, after a global discount sweep over every cell.

\subsection{Dual Gaussian posteriors}
\label{sec:method-posteriors}

For each arm $a \in \{1,\ldots,K\}$ and bucket $k \in \{1,\ldots,K_b\}$ (where $K_b$ is the number of context buckets, default $4$, one per task type), we maintain two posteriors:
\begin{itemize}
\item \textbf{Q-posterior}, tracking the per-arm utility $\ut = q_t - \nu f_t$. Conjugate Normal-Normal with prior $\mathcal{N}(\mu_{0,q}, \sigma^2_{0,q})$ and likelihood variance $\sigma^2_q$.
\item \textbf{C-posterior}, tracking the per-arm raw cost $c_t$. Same Normal-Normal form, with prior $\mathcal{N}(\bar c_a, \sigma^2_{0,c})$ where $\bar c_a$ is the spec base price for arm $a$, and likelihood variance $\sigma^2_c$.
\end{itemize}

Each cell carries the standard discounted sufficient statistics $n^{a,k}_t \in \mathbb{R}_{\geq 0}$ (effective sample count) and $S^{a,k}_t \in \mathbb{R}$ (weighted sum of observations), with the analogous pair for the cost posterior. Posterior parameters follow the closed-form Normal-Normal expressions, e.g.\ for Q:
\begin{equation}
\label{eq:posterior-q}
\Sigmaq{a}{k} = \left(\frac{1}{\sigma^2_{0,q}} + \frac{n^{a,k}_t}{\sigma^2_q}\right)^{-1}, \qquad
\thetaq{a}{k} = \Sigmaq{a}{k} \left(\frac{\mu_{0,q}}{\sigma^2_{0,q}} + \frac{S^{a,k}_t}{\sigma^2_q}\right).
\end{equation}
The C-posterior uses the analogous form with $(\thetac{a}{k}, \Sigmac{a}{k})$, $\bar c_a$ in place of $\mu_{0,q}$, $\sigma^2_c$ in place of $\sigma^2_q$, and a separate pair of sufficient statistics over realized costs.

\paragraph{Why a learned cost posterior.}
In stationary settings with known prices, a fixed scalar $\bar c_a$ would suffice for cost. The agent-native regime drops both assumptions: prices can change mid-run (S3 in our benchmark), and the agent only sees prices through quoted observations from the market. Treating cost as a Gaussian posterior over realized $c_t$ lets the policy detect price shocks the same way the Q-posterior detects quality shocks---through the effective sample count $n^{a,k}_t$ and the discount that reweights stale observations. We demonstrate empirically (\S\ref{sec:eval-ablations}) that disabling this and pinning cost to $\bar c_a$ collapses adaptation to S3.

\subsection{Discounted sufficient statistics}
\label{sec:method-discount}

We apply a multiplicative discount $\gamma \in (0, 1]$ to every cell's $n$ and $S$ at the start of each round, before any new observation arrives:
\begin{equation}
n^{a,k}_t = \gamma \cdot n^{a,k}_{t-1}, \qquad
S^{a,k}_t = \gamma \cdot S^{a,k}_{t-1}.
\end{equation}
After the round's chosen arm produces an observation, only its bucket increments:
\begin{equation}
n^{a_t, k_t}_t \mathrel{+}= 1, \quad
S^{a_t, k_t}_t \mathrel{+}= \ut\,\,(\text{Q-posterior}), \quad
{S_c}^{a_t, k_t}_t \mathrel{+}= c_t\,\,(\text{C-posterior}).
\end{equation}
Discounting applies regardless of whether a cell was pulled this round---non-stationarity is wall-clock, not pull-conditional. We use $\gamma = 0.999$ for both posteriors (effective half-life $\approx 693$ rounds), chosen to react cleanly within the shock windows we evaluate (S2 covers $2{,}500$ rounds, S3 has $9{,}000$ post-shock rounds) without forgetting the stationary baseline too quickly. Sensitivity to $\gamma$ is in [/cc TODO sensitivity appendix].

\subsection{Decision rule}
\label{sec:method-decision}

At each round, the wallet exposes $\lambda_t$ and the affordable set $\mathcal{A}_t$. The policy reads $\lamn$ from the wallet, samples both posteriors per affordable arm, and ranks them:
\begin{equation}
\label{eq:padct-sample}
\hatu{a}{k_t} \sim \mathcal{N}\!\big(\thetaq{a}{k_t},\, \Sigmaq{a}{k_t}\big), \qquad
\hatc{a}{k_t} \sim \mathcal{N}\!\big(\thetac{a}{k_t},\, \Sigmac{a}{k_t}\big),
\end{equation}
\begin{equation}
\label{eq:padct-score}
\tilde{\hat c}^{a,k_t} = \mathrm{clip}\!\big(\hatc{a}{k_t}/c_{\max},\, 0,\, 1\big), \qquad
\pi_a = (1 - \lamn)\,\hatu{a}{k_t} - \lamn\,\tilde{\hat c}^{a,k_t}.
\end{equation}
The chosen arm is $a_t = \arg\max_{a \in \mathcal{A}_t} \pi_a$. Algorithm~\ref{alg:padct} fixes the per-round procedure end-to-end, with each block annotated by the PA-DCT component it implements.

\begin{algorithm}[t]
\caption{PA-DCT (Payment-Aware Discounted Contextual Thompson sampling). Block annotations P, D, C, TS mark the four ablatable components.}
\label{alg:padct}
\begin{algorithmic}[1]
\Require Provider set $\{1,\ldots,K\}$ with spec costs $\bar c_a$; discount $\gamma$; failure penalty $\nu$; cost normalizer $c_{\max}$; encoder $k(\cdot)$; priors $(\mu_{0,q}, \sigma^2_{0,q}, \sigma^2_q)$ and $(\bar c_a, \sigma^2_{0,c}, \sigma^2_c)$; wallet sharpness $\alpha$.
\For{$(a, k) \in \{1,\ldots,K\} \times \{1,\ldots,K_b\}$}
  \State init Q-posterior$_{a,k} \leftarrow \mathcal{N}(\mu_{0,q}, \sigma^2_{0,q})$;\quad C-posterior$_{a,k} \leftarrow \mathcal{N}(\bar c_a, \sigma^2_{0,c})$
\EndFor
\For{$t = 0, 1, \ldots, T-1$}
  \State \emph{(D)} discount every Q- and C-cell by $\gamma$
  \State observe $\mathbf{x}_t$, $B_t$;\quad $k_t \leftarrow k(\mathbf{x}_t)$
  \State \emph{(P)} $\lambda_t \leftarrow \exp(\alpha \cdot \bex_t)$;\quad $\lamn \leftarrow \lambda_t / (1 + \lambda_t)$
  \State $\mathcal{A}_t \leftarrow \{a : B_t \geq \mathrm{effective\_price}(a, t)\}$
  \If{$\mathcal{A}_t = \emptyset$} halt (bankruptcy) \EndIf
  \For{$a \in \mathcal{A}_t$}
    \State \emph{(TS, C)} sample $\hatu{a}{k_t} \sim$ Q-posterior$_{a, k_t}$;\quad $\hatc{a}{k_t} \sim$ C-posterior$_{a, k_t}$
    \State $\tilde{\hat c}^{a,k_t} \leftarrow \mathrm{clip}(\hatc{a}{k_t}/c_{\max}, 0, 1)$;\quad $\pi_a \leftarrow (1-\lamn)\,\hatu{a}{k_t} - \lamn\,\tilde{\hat c}^{a,k_t}$
  \EndFor
  \State $a_t \leftarrow \arg\max_{a \in \mathcal{A}_t} \pi_a$
  \State pay via the protocol; observe $(q_t, c_t, f_t)$;\quad $\ut \leftarrow q_t - \nu f_t$
  \State update Q-posterior$_{a_t, k_t}$ with $\ut$;\quad update C-posterior$_{a_t, k_t}$ with $c_t$
  \State wallet.\texttt{record\_spend}$(c_t)$
\EndFor
\end{algorithmic}
\end{algorithm}

\subsection{Hyperparameters}
\label{sec:method-hyperparams}

We fix the following constants across every experiment in the paper:
\begin{itemize}
\item Quality posterior: $\mu_{0,q} = 0.5$, $\sigma^2_{0,q} = 1.0$, $\sigma^2_q = 0.09$ (allowing substantial heterogeneity in observed quality).
\item Cost posterior: $\sigma^2_{0,c} = 10^{-4}$, $\sigma^2_c = 10^{-6}$ (tight, because realized cost is nearly deterministic given the provider and scenario).
\item Discounts: $\gamma_q = \gamma_c = 0.999$.
\item Reward weights and bandit constants: $\nu = 0.5$, $\alpha = 2$, $c_{\max} = \$0.01$.
\item Bucket count: $K_b = 4$ (one per task type, see \S\ref{sec:eval-setup}).
\end{itemize}
None of these are tuned per scenario or per seed.

\subsection{Theoretical properties}
\label{sec:method-theory}

We establish four properties of PA-DCT: reward boundedness (Prop.~\ref{prop:reward-bounded}), posterior consistency in the stationary regime (Prop.~\ref{prop:consistency}), a Bayesian regret bound in the stationary regime (Thm.~\ref{thm:stationary-regret}), and a per-arm adaptation rate after an abrupt change-point (Thm.~\ref{thm:adaptation}). Proof sketches are given here; full proofs are deferred to Appendix~\ref{appendix:proofs}.

\paragraph{Setting and assumptions.}
Throughout this subsection we assume:
\begin{itemize}
\item (A1) per-cell stationary or piecewise-stationary distribution of $(q_t, c_t, f_t)$ given the chosen arm and bucket;
\item (A2) sub-Gaussian noise on the utility $u_t = q_t - \nu f_t$ and on the realized cost $c_t$;
\item (A3) bounded means $|\mu_u^{a,k}| \leq 1$ and $\mu_c^{a,k} \in [0, c_{\max}]$ for every arm $a$ and bucket $k$.
\end{itemize}
Theorem~\ref{thm:stationary-regret} (the regret reduction) further requires two strengthenings used to invoke standard contextual Thompson-sampling bounds:
\begin{itemize}
\item (A2$^\prime$) \emph{Well-specified Gaussian Bayesian model}: the policy's likelihood proxy variances $\sigma^2_q, \sigma^2_c$ match the true noise variances $\sigma_u^2, \sigma_{c,\mathrm{true}}^2$, and the per-cell prior is the Bayesian-correct prior over $(\mu_u^{a,k}, \mu_c^{a,k})$. Misspecification (proxy $\neq$ true) is discussed in \S\ref{app:proof-limitations}.
\item (A3$^\prime$) \emph{Cost margin}: $\mu_c^{a,k} \in [\delta_c, c_{\max} - \delta_c]$ for some uniform margin $\delta_c > 0$, so the cost-clip event $\{\hat c^{a,k}/c_{\max} \notin [0,1]\}$ has probability vanishing exponentially in the cost-posterior precision (Lemma in Appendix~\ref{app:proof-stationary}).
\end{itemize}
Propositions~\ref{prop:reward-bounded}--\ref{prop:consistency} and Theorem~\ref{thm:adaptation} use only (A1)--(A3); only Theorem~\ref{thm:stationary-regret} requires (A2$^\prime$) and (A3$^\prime$).

\begin{proposition}[Reward boundedness]
\label{prop:reward-bounded}
For every round $t$, the payment-aware reward satisfies $\rt \in [-1,+1]$.
\end{proposition}

\begin{proof}[Proof sketch]
$\ut = q_t - \nu f_t$ with $q_t \in [0,1]$, $f_t \in \{0,1\}$, and $\nu = 0.5$ gives $|\ut| \leq 1$. Combined with $\ctilde \in [0,1]$ and $\lamn \in [0,1]$, the triangle inequality on $\rt = (1-\lamn)\ut - \lamn\,\ctilde$ gives $|\rt| \leq (1-\lamn) + \lamn = 1$. See Appendix~\ref{app:proof-bounded}.
\end{proof}

\begin{proposition}[Posterior consistency, stationary regime]
\label{prop:consistency}
Fix $(a,k)$ and assume the utility distribution at $(a,k)$ is stationary with finite variance and mean $\mu_u^{a,k}$. With no discount ($\gamma=1$), the Q-posterior mean $\thetaq{a}{k}$ converges almost surely to $\mu_u^{a,k}$ as the cell's pull count $n^{a,k} \to \infty$. The same holds for the C-posterior and $\mu_c^{a,k}$.
\end{proposition}

\begin{proof}[Proof sketch]
Standard Normal-Normal conjugacy gives a posterior mean of the form $(\mu_0/\sigma_0^2 + S/\sigma^2)/(1/\sigma_0^2 + n/\sigma^2)$. As $n\to\infty$ the prior term becomes negligible and $S/n \to \mu$ by the strong law of large numbers. See Appendix~\ref{app:proof-consistency}.
\end{proof}

\paragraph{Effective sample size under discount.}
A key quantity for the analysis under $\gamma < 1$ is the discounted effective sample size $n^{a,k}_t \leq (1-\gamma^{t+1})/(1-\gamma) \leq 1/(1-\gamma)$ (Lemma~\ref{lem:ess} in Appendix~\ref{appendix:proofs}). For $\gamma = 0.999$ this caps the effective horizon at roughly $1{,}000$ rounds, defining the natural responsiveness of the policy to drift.

\begin{theorem}[Regret reduction to contextual Gaussian Thompson sampling]
\label{thm:stationary-regret}
Under (A1), (A2$^\prime$), (A3$^\prime$), and $\gamma = 1$, define the \emph{policy-state pseudo-regret} of PA-DCT against the wallet-coupled oracle as
\begin{equation}
\label{eq:policy-state-regret}
\Reg_T \;:=\; \sum_{t=0}^{T-1} \Big[(1-\lamn)\,\mu_u^{a^\star_t, k_t} - \lamn\,\mu_c^{a^\star_t, k_t}/c_{\max}\Big] - \Big[(1-\lamn)\,\mu_u^{a_t, k_t} - \lamn\,\mu_c^{a_t, k_t}/c_{\max}\Big],
\end{equation}
where $a^\star_t \in \arg\max_{a \in \mathcal{A}_t}$ of the bracketed quantity at the policy's wallet pressure $\lamn$, and $a_t$ is PA-DCT's chosen arm. Then PA-DCT reduces, conditional on the bucket-visit sequence, to per-cell K-armed Gaussian Thompson sampling with a one-dimensional time-varying linear-context payoff. Under (A2$^\prime$) the per-cell Bayesian regret obeys the standard bound
\begin{equation}
\label{eq:cell-regret}
\mathbb{E}\!\left[\Reg^{\mathrm{cell}\,k}_{T_k}\right] \;\leq\; C_0(\sigma_q^2, \sigma_c^2, \sigma^2_{0,q}, \sigma^2_{0,c}, c_{\max})\cdot \sqrt{K\, T_k\, \log T_k},
\end{equation}
inherited from \citet{russo2014learning} Theorem~1 (and \citealt{lattimore2020bandit} Ch.~36 for the K-armed Gaussian-prior version). Aggregating across $K_b$ buckets via $\sum_k T_k = T$ and Cauchy--Schwarz,
\begin{equation}
\label{eq:stationary-regret}
\mathbb{E}\!\left[\Reg_T\right] \;\leq\; C \cdot \sqrt{K\,K_b\,T \log T},
\end{equation}
for a constant $C$ depending on the prior and likelihood variances and the cost margin $\delta_c$ of (A3$^\prime$).
\end{theorem}

\begin{proof}[Proof sketch]
Three steps. (i) Sampling $\hat u^{a,k}$ and $\hat c^{a,k}$ independently is equivalent in distribution to sampling once from the Gaussian posterior on the linear combination $\pi_a := (1-\lamn)\hat u^{a,k} - \lamn\,\hat c^{a,k}/c_{\max}$ (same mean, variance $(1-\lamn)^2 \Sigmaq{a}{k} + \lamn^2 \Sigmac{a}{k}/c_{\max}^2$). (ii) The wallet weight $\lamn$ enters as part of the time-varying \emph{context} and is observed before sampling, so within bucket $k$ the policy is contextual K-armed Gaussian TS with feature $\phi_t = (1-\lamn, -\lamn/c_{\max}) \in \mathbb{R}^2$, and the comparator in (\ref{eq:policy-state-regret}) uses the same $\phi_t$. The clip $\mathrm{clip}(\hat c/c_{\max},0,1)$ is handled via Lemma~\ref{lem:clip} in Appendix~\ref{app:proof-stationary}: under (A3$^\prime$), the per-round event $\{\hat c^{a,k}/c_{\max} \notin [0,1]\}$ has probability decaying exponentially in the cost-posterior precision and contributes only a lower-order $O(\log T)$ term to the regret. (iii) Russo--Van Roy 2014 Theorem~1 applied per cell gives~(\ref{eq:cell-regret}), and Cauchy--Schwarz aggregates to~(\ref{eq:stationary-regret}). Full proof in Appendix~\ref{app:proof-stationary}.
\end{proof}

\begin{remark}[Relation to the empirical \emph{True Oracle} of \S\ref{sec:eval-comparators}]
The wallet-aware oracle in Thm.~\ref{thm:stationary-regret} uses the \emph{policy's} wallet pressure $\lamn$ at each round (i.e., it answers ``what is the best arm given the current state of the policy's wallet?''). The \emph{True Oracle} reported in \S\ref{sec:eval-comparators} is a different comparator: it owns its own wallet and so its own $\lamn$ trajectory, hence achieves at least as much reward as the wallet-aware oracle. Theorem~\ref{thm:stationary-regret} therefore upper-bounds a regret quantity that is no larger than the empirical regret reported in \S\ref{sec:eval-main}; we did not attempt to bound the True-Oracle regret directly because doing so requires controlling the divergence between the policy's and the oracle's wallet trajectories.
\end{remark}

\begin{theorem}[Adaptation rate under an abrupt change, exact form]
\label{thm:adaptation}
Suppose the true utility mean of arm $a$ in bucket $k$ shifts between round $t^\star-1$ and round $t^\star$, so that $t^\star$ is the first post-shock round; the pre-shock mean is $\mu_{\mathrm{old}}$, the post-shock mean is $\mu_{\mathrm{new}}$, and $\Delta = |\mu_{\mathrm{old}} - \mu_{\mathrm{new}}|$. Define
\[
A := \gamma^{t-t^\star+1}\, n^{a,k}_{t^\star-1}, \qquad B := \sum_{i=1}^{m} \gamma^{t-s_i}, \qquad \kappa_q := \sigma^2_q / \sigma^2_{0,q},
\]
where $\{s_1 < \cdots < s_m\} \subseteq [t^\star, t]$ are the post-shock pulls of $(a,k)$. Then the bias of the Bayesian Q-posterior mean of PA-DCT satisfies the \emph{exact} bound
\begin{equation}
\label{eq:adaptation}
\big|\,\mathbb{E}[\theta^{a,k}_{q,t} \mid m,\, s_1,\ldots,s_m] - \mu_{\mathrm{new}}\,\big| \;\leq\; \frac{A\,\Delta \;+\; \kappa_q\,|\mu_{0,q} - \mu_{\mathrm{new}}|}{\kappa_q + A + B}.
\end{equation}
\end{theorem}

\begin{proof}[Proof sketch]
The discounted Normal-Normal posterior mean is $\theta^{a,k}_{q,t} = (\kappa_q\,\mu_{0,q} + S^{a,k}_t)/(\kappa_q + n^{a,k}_t)$, after multiplying numerator and denominator of (\ref{eq:posterior-q}) by $\sigma^2_q$. Decomposing $S^{a,k}_t$ and $n^{a,k}_t$ into pre- and post-shock contributions (Lemma~\ref{lem:ess} in Appendix~\ref{appendix:proofs} for the contribution rule) and substituting the in-expectation values yields~(\ref{eq:adaptation}) by triangle inequality on the numerator. See Appendix~\ref{app:proof-adaptation}.
\end{proof}

\begin{remark}[Three regimes]
\label{rem:adapt-regimes}
The exact bound~(\ref{eq:adaptation}) admits three useful simplifications:
\begin{itemize}
\item \textbf{Prior-dominated} ($\kappa_q \gg A + B$, i.e., very few pulls): bias $\to |\mu_{0,q} - \mu_{\mathrm{new}}|$. The posterior is uninformative, and the prior controls the bias.
\item \textbf{Data-dominated} ($\kappa_q \ll A + B$): bias $\to \Delta\, A/(A+B)$, which is the term we discount through additional pulls.
\item \textbf{Worst case in $\Delta$}: even when $\kappa_q$ is non-negligible, bias $\leq \Delta + \kappa_q\,|\mu_{0,q} - \mu_{\mathrm{new}}|/(\kappa_q + A + B)$, separating the data error from the irreducible prior error.
\end{itemize}
\end{remark}

\begin{corollary}[Recovery horizon, data-dominated regime]
\label{cor:recovery}
Suppose $(a,k)$ is pulled at every post-shock round, so $m = \tau$ with $\tau := t - t^\star + 1$, and assume the data-dominated regime $\kappa_q \ll A + B$ (well approximated whenever $\tau \gtrsim \kappa_q$, since $A + B \geq B \geq 1$). Then the bound~(\ref{eq:adaptation}) reduces to
\[
\big|\,\mathbb{E}[\theta^{a,k}_{q,t}] - \mu_{\mathrm{new}}\,\big| \;\lesssim\; \Delta\, \gamma^{\tau} \;+\; O\!\big(\kappa_q / (A + B)\big),
\]
where the data-dominated term $\Delta \gamma^\tau$ shrinks geometrically in $\tau$ and the residual prior term is $O(\kappa_q (1-\gamma))$ in the asymptotic regime $A + B \to 1/(1-\gamma)$. Hence the recovery condition $\Delta \gamma^\tau \leq \delta$ requires
\[
\tau \;\geq\; \frac{\log(\Delta/\delta)}{\log(1/\gamma)} \;=\; O\!\big(\log(\Delta/\delta)/(1-\gamma)\big).
\]
\end{corollary}

\begin{proof}[Proof sketch]
Plug $A \leq \gamma^\tau/(1-\gamma)$ (Lemma~\ref{lem:ess}) and $B = (1-\gamma^\tau)/(1-\gamma)$ into~(\ref{eq:adaptation}); the data term gives $\gamma^\tau$ exactly and the prior term is bounded by $\kappa_q\,|\mu_{0,q}-\mu_{\mathrm{new}}|/(\kappa_q + 1/(1-\gamma)) = O(\kappa_q (1-\gamma))$ once $\kappa_q$ is dwarfed by the post-shock evidence mass. The hyperparameters of \S\ref{sec:method-hyperparams} ($\kappa_q = 0.09$, $\gamma = 0.999$) place the prior term well below the empirical noise floor for any moderate $\tau$. See Appendix~\ref{app:proof-adaptation}.
\end{proof}

\paragraph{What this analysis does and does not give us.}
Together: PA-DCT operates on a bounded reward (Prop.~\ref{prop:reward-bounded}), its beliefs become correct in the limit (Prop.~\ref{prop:consistency}), it inherits standard contextual Gaussian-TS regret bounds in the stationary well-specified regime (Thm.~\ref{thm:stationary-regret}), and its discount factor $\gamma$ controls an explicit per-arm recovery horizon after an abrupt change (Thm.~\ref{thm:adaptation}, Cor.~\ref{cor:recovery}). What this analysis does \emph{not} give: (i) a tight $C$ in~(\ref{eq:stationary-regret}); the constant inherits its tightness from \citet{russo2014learning} and we do not claim to improve it. (ii) a regret bound under \emph{misspecified} likelihoods (relaxing (A2$^\prime$); see App.~\ref{app:proof-limitations}). (iii) a coupled regret bound for piecewise-stationary regimes with budget pressure---combining contextual TS, discount-induced drift adaptation, and the wallet-pressure feedback that ties the policy's $\lamn$ trajectory to its own past actions; this would require a regret bound of the shape $\tilde O(\sqrt{K K_b T} + M \Delta_{\max}/(1-\gamma))$ in the spirit of \citet{garivier2011discounted,trovo2020sliding} and does not follow from the cited prior results. (iv) a True-Oracle regret bound (relaxing the wallet-coupled comparator); see Remark immediately following Thm.~\ref{thm:stationary-regret} and App.~\ref{app:proof-limitations}. Bandits-with-knapsacks bounds~\citep{badanidiyuru2018bandits} are adjacent but do not transfer cleanly because our budget constraint is soft. We report empirical regret against the True Oracle in \S\ref{sec:eval-main}.
